\subsection{An application of the existing optimization theorem}
\label{sec:algorithm-application}

The full-support theorem also verifies the hypothesis of an existing
algorithmic result. We use the operation count of El Alaoui, Montanari,
and Sellke~\cite[Remark~2.1, Theorem~3, and Corollary~2.2 of the cited
preprint]{EAMS2021}: real sums and products have unit cost, and the scalar
functions and normalization constants are those prescribed in their
algorithm for the chosen accuracy. This is an instance-size bound;
the argument below supplies no separate complexity estimate for numerical
preprocessing of those functions or for finite-precision arithmetic.

\begin{corollary}[Asymptotically optimal SK optimization]
\label{cor:algorithm}
For every fixed $\epsilon>0$, there is a randomized algorithm with
operation count at most $C(\epsilon)n^2$ in the model just described,
where $C(\epsilon)$ is independent of $n$, that takes the Gaussian
couplings of $H_n$ as input and returns
$\sigma^{\rm alg}\in\{-1,1\}^n$ such that
\[
 \Pr\left\{
 H_n(\sigma^{\rm alg})\geq
 \max_{\sigma\in\{-1,1\}^n}H_n(\sigma)-\epsilon n
 \right\}\longrightarrow1
 \qquad(n\to\infty).
\]
The probability includes the couplings and the algorithm's randomness.
\end{corollary}

\begin{proof}
Theorem~\ref{thm:full-support} gives a strictly increasing minimizer
in $\mathcal U$, exactly the no-overlap-gap hypothesis in
\cite[Assumption~2]{EAMS2021} for $\xi(t)=t^2/2$.
We apply their Corollary~2.2, with its operation count from Theorem~3.

To match their Gaussian quadratic input, form a symmetric matrix $A$
with $A_{ij}=J_{ij}/\sqrt n$ for $i<j$ and independent diagonal entries
$A_{ii}=\sqrt{2/n}\,G_i$, where the $G_i$ are independent standard
Gaussians. The quadratic extension
$\widehat H_n(x)=\frac12x^{\mathsf T}Ax$ has covariance
$(2n)^{-1}(\sigma^{\mathsf T}\tau)^2$ on the hypercube, and
\[
 \widehat H_n(\sigma)
   =H_n(\sigma)+\frac{1}{\sqrt{2n}}\sum_{i=1}^nG_i.
\]
The second term is independent of $\sigma$, so every optimization gap
is identical for $H_n$ and $\widehat H_n$.

Evaluating $\nabla\widehat H_n(x)=Ax$ costs $O(n^2)$ operations.
The part with distinct indices is $H_n(x)$, and evaluating one coordinate
of its gradient costs $O(n)$. Thus the bound
$C(\epsilon)(\chi+n)+n\chi_1$ in the cited theorem is
$C(\epsilon)n^2$, after changing the constant. Its probability guarantee
and the equality of the optimization gaps give the assertion.
\end{proof}

This corollary invokes the finite-disorder message-passing theorem of
El Alaoui, Montanari, and Sellke. Proposition~\ref{prop:gaussian} supplies
the separate scalar approximation proved above.
